\textbf{Problems}

\begin{enumerate}[label=\thesection.\arabic*.,ref=\thesection.\arabic*,start=12,series=book-problems]

\item\label{ex:closed-graph-implies-bounded-inverse} 证明闭图定理蕴含有界逆定理.

\item\label{ex:orthogonal-projection-open-map} 设 $M$ 是 Hilbert 空间 $H$ 的闭子空间. 不使用开映射定理, 证明正交投影算子 $P:H\to M$ 将 $H$ 中的开集映为 $M$ 中的开集.

\item\label{ex:surjective-hilbert-operator-null-complement-bijection} 设 $A:H\to Y$ 是从 Hilbert 空间 $H$ 到 Banach 空间 $Y$ 的有界满射线性算子. 设 $\widetilde A:\mathcal{N}(A)^\perp\to Y$ 为 $A$ 在定义域 $\mathcal{N}(A)^\perp$ 上的限制. 证明 $\widetilde A$ 是连续双射.

\item\label{ex:open-map-need-not-preserve-closed-sets} 若映射将每个开集映为开集, 则称它是开映射. 证明开映射未必将闭集映为闭集.

\item\label{ex:incomplete-finite-sequence-bounded-inverse-counterexample} 设 $X$ 为仅有有限个非零项的序列 $x=\{x_1,x_2,x_3,\ldots\}$ 所组成的空间, 范数为\[\|x\|:=\sup_{i\in\mathbb{N}}|x_i|.\]用下式定义 $T:X\to X$:\[Tx:=\left\{x_1,\dfrac{x_2}{2},\dfrac{x_3}{3},\ldots\right\}.\]证明 $T$ 线性且有界, 但 $T^{-1}$ 无界. 为什么这不与有界逆定理矛盾?

\end{enumerate}
